\bookchapter{第十二章：音声·映像エンターテインメントと読書}

\chapterintro{一、音楽鑑賞：どのAppを選ぶか、どう聴くと快適か}

\booksection{音楽Appは現在機能がほぼ同じですが、それぞれ強みがあり、選び方次第でかなり快適になります：}

網易雲音楽

\begin{itemize}
  \item \textbf{：強みはプレイリストコミュニティとコメント欄。「雨の日に聴く曲」や「勉強疲れ時にリラックスできるインスト」を検索すれば、他人が整理したプレイリストがすぐ見つかります。コメント欄が曲より面白いことも。欠点は一部の楽曲の著作権が不完全で、聴いている途中でグレーになる曲があること。}QQ音楽
  \item \textbf{：著作権が最も充実、ジェイ·チョウ、五月天などの人気アーティストの楽曲は基本的にQQ音楽にしかありません。音質オプションも豊富で、会員はハイレゾ音質を利用可能。「聴きたい曲がすべて検索できる」ことを重視するなら、QQ音楽が最も確実。}Apple Music
  \item \textbf{：Appleユーザーのファーストチョイス。インターフェースがクリーン、広告なし、音質が良く、ロスレスやDolby Atmosに対応。ただしレコメンドアルゴリズムは国内Appほど中国語ユーザーの好みを理解せず、プレイリスト文化も網易雲ほど発展していません。シンプルな体験を好み、ソーシャル機能を気にしない人に適しています。}Spotify
  \item \textbf{：世界的にユーザーが最も多く、レコメンドアルゴリズムは公認トップ。「毎週の発見」と「デイリーミックス」の精度が驚異的で、長く使うほどユーザーを理解。欠点は無料版の広告が多く、一部の中国語曲の著作権カバーが不完全。利用には特殊なネットワークが必要。}聴曲の小技
\end{itemize}

\textbf{：Bluetoothイヤホンを使っている場合、音楽Appの設定で「音響調整」や「イコライザー」があるか確認し、低音と高音を少し調整すると効果が明確に向上。また、ストリーミング聴曲は通信量を消費するので、設定で「Wi-Fi時のみロスレス再生」をオンにするか、お気に入りのプレイリストを事前ダウンロードしてオフラインで聴きましょう。}二、ポッドキャスト（Podcast）：通勤路上の「耳の伴侶」

\booksection{ポッドキャストは、聴いたことがない人は自分と無関係だと思い、一度聴き始めると戻れなくなります。本質は「ラジオ番組のNetflix版」——科技を語る人、歴史を語る人、物語を読む人、どんなジャンルでも揃っています。}

小宇宙アプリ

\begin{itemize}
  \item \textbf{：国内で最も良くできたポッドキャストApp、インターフェースが清潔で余計な機能なし。「毎日のおすすめ」と「編集者厳選」が新番組の発見を助け、コメント欄も活発。私が最もよく聴くのは《忽左忽右》（文化歴史系）と《不合時宜》（社会問題系）で、毎回賢い人たちの会話を傍聴しているようです。}Appleポッドキャスト
  \item \textbf{：iPhoneに標準搭載、追加App不要。CarPlayとの統合が最も良く、運転中に便利。}聴取のコツ
  \item \textbf{：ほとんどのポッドキャストAppは1.2倍速や1.5倍速再生に対応、慣れると時間を大幅に節約。私は中国語は1.2倍、英語は1.5倍で聴いています。また「新エピソードの自動ダウンロード」と「再生完了後の自動削除」を必ずオンに。そうしないとスマホに聴き終わった音声ファイルが溜まっていきます。}三、動画視聴：プレーヤーからストリーミングまでの正しい開き方
\end{itemize}

\booksection{多くの人が知りませんが、スマホで再生できる形式はストリーミングAppが提供するものよりはるかに多い。良いローカルプレーヤーで多くの問題が解決します。}

VLCプレーヤー

\begin{itemize}
  \item \textbf{：オープンソース無料、ほぼあらゆる形式を再生——MKV、AVI、FLAC、さらにはISOディスクイメージも。パソコンでは再生できるがスマホでは再生できない動画をダウンロードした場合、十中八九フォーマット問題で、VLCを入れれば解決。外付け字幕にも対応、.srt字幕ファイルを動画と同じフォルダに置けばVLCが自動的に読み込みます。}nPlayer
  \item \textbf{：より洗練された体験を求めるなら、この有料Appは買い。インターフェースが美しく、ジェスチャー操作が快適——左側上下スワイプで明るさ調整、右側上下スワイプで音量調整、左右スワイプで早送り巻き戻し。DTSやDolby音響にも対応、映画視聴時の臨場感が大幅に向上。}Bilibili
  \item \textbf{：もはや二次元サイトだけでなく、学習リソースが異常なほど豊富。大学入試、プログラミング、撮影、料理、あらゆるチュートリアルが揃います。弾幕文化はBilibiliの魂で、動画視聴時に弾幕をオンにすると、一緒に見ている人がいる感覚。ただし学習動画視聴時は集中力を削がれるので弾幕オフ推奨。}テレビにキャスト
  \item \textbf{：スマホ画面は大きくてもテレビには敵いません。同じWi-Fi環境で、動画App右上の「TV」キャストボタンをタップすれば、テレビに画面を投影可能。ほとんどのスマートテレビやプロジェクターが対応、ドラマや映画の視聴体験が大幅に向上。iPhoneユーザーはAirPlayで画面全体をミラーリング可能。}四、電子書籍読書：スマホでもしっかり読める
\end{itemize}

\booksection{スマホでの読書は目に悪い、集中力を削ぐと言われますが、確かに一理あります。しかしスマホ読書の利点も明確——いつでもどこでも、細切れ時間を最大化活用。私は地下鉄、行列、寝る前にスマホで多くの本を読みました。}

WeChat読書

\begin{itemize}
  \item \textbf{：現在国内で体験が最も良い読書App。書庫が充実、出版書もネット小説も揃います。「読書時間でコイン交換」モデルで、多く読めば無料で本と交換でき、学生に優しい。「アンダーライン」機能で他人がどの段落にマークしたか見え、他人の注釈した良い文に共感することも。ただし「読書ランキング」は少し競争が激しく、「時間稼ぎのために読む」になりがちなので、オフ推奨。}Kindle App
  \item \textbf{：Kindleデバイスを持っていれば、スマホ版Kindle Appで読書進捗を同期。Kindleの正規電子書籍はレイアウト品質が海賊版PDFよりはるかに良いですが、Kindle中国ストアは閉鎖され、海外リージョンへの切り替えが必要。}Apple Books（ブックス）
  \item \textbf{：iPhone標準、EPUBとPDFの読書に適しています。インターフェースが極めてシンプル、ページめくりアニメーションが快適。ネットから電子書籍ファイルをダウンロードしたら、「ブックスに共有」でインポート可能。}読書設定の提案
  \item \textbf{：どのAppでも、背景を暖黄色調（アイケアモード）にし、フォントを少し大きくし、行間を広めに。夜の読書では「ダークモード」をオン——黒背景白文字で目が楽。絶対に白背景黒文字小フォントを暗所で長時間読まないで、目の疲れは蓄積します。}五、スマホゲーム：楽しく、スマホにも優しく
\end{itemize}

\booksection{スマホゲームへの私の態度は：遊んで良いが、逆にコントロールされないように。ゲームツールを合理的に使えば、より快適に遊べ、本業も疎かになりません。}

ゲームモード/ゲームスペース

\begin{itemize}
  \item \textbf{：最近の主流スマホにはこの機能があります。オンにするとスマホが自動的にバックグラウンドクリア、通知ブロック、誤操作防止——団体戦中に突然電話が鳴るような事態をゲームモードが防いでくれます。一部のスマホはパフォーマンスモード/E-sportsモードにも対応、プロセッサを強制フル稼働させますが、代償は発熱と消費電力の増大。}画質とフレームレートの調整
  \item \textbf{：無闇に最大にしない。まずフレームレート（滑らかさ）を確保してから画質を追求。60fpsは30fpsよりはるかに快適ですが、120fpsはスマホに負荷が大きく、ミドルレンジ機では無理する必要なし。競技系ゲーム（王者、PUBGなど）では、フレームレートの優先度は画質よりはるかに高い——滑らかな操作感は綺麗な芝生より百倍重要。}ゲーム中のスマホ発熱対策
  \item \textbf{：これは正常ですが、長期の高温はバッテリー寿命を損ないます。いくつかの冷却策：スマホケースを外して放熱、画面輝度を70\%程度に下げる（輝度は高発熱源）、充電しながらゲームをしない（発熱が重なり、バッテリーが急速に劣化）。スマホが握れないほど熱い場合は、すぐにゲームを終了して休ませましょう。}依存防止ツール
  \item \textbf{：iOSの「スクリーンタイム」とAndroidの「デジタルウェルビーイング」で、ゲームAppに1日の時間制限を設定可能。私は平日最大1時間、週末2時間に設定。時間に達するとAppアイコンがグレーアウトし、タップできなくなります。この制限は本当に効果的で、形だけのものではありません。}クラウドゲーム
  \item \textbf{：スマホ性能が足りないが大作を遊びたい場合は、クラウドゲームプラットフォーム（Tencent STARTクラウドゲーム、NetEaseクラウドゲームなど）を試す。ゲームはリモートサーバーで実行され、映像が動画ストリームとしてスマホに送信。ネットワーク要件は高いがスマホ性能はほぼ不要。欠点は遅延が避けられないことで、戦略系、RPG系に適し、FPS系には不向き。}やってみよう：スマホエンターテインメントシステムを構築
\end{itemize}

\begin{tcbpractice}{音楽Appでお気に入りのアルバムやプレイリスト3つをダウンロードし、オフライン再生を確保。}
\begin{enumerate}
  \item 小宇宙またはAppleポッドキャストで興味あるチャンネル3つを登録し、1エピソードをダウンロードして試聴。
  \item VLCプレーヤーをインストールし、パソコンからコピーした動画ファイルの再生を試す。
  \item WeChat読書で読みたい本を探し、暖色背景とアイケアモードを設定、自分に合ったフォントサイズに調整。
  \item ゲームをする場合、ゲームモードを開いて最適化設定を確認（迷惑防止、フレームレート優先など）。
  \item 第十三章：ツール系Appのおすすめ
\end{enumerate}

\end{tcbpractice}
